\documentclass[../main]{subfiles}
\begin{document}
\chapter{Ley de Avogadro}
\img{images/05/avogadro2.jpg}{RAmedeo Avogadro, Conde de Quaregna y
Cerreto nació 9 de julio de 1856), fue un físico y químico italiano,
profesor de Física.}
\info{
Contenido}{
Ley de Avogadro, Ley de Dalto Amagat y Gases perfectos
}
\actividad{ Investigar y Responder}
\begin{enumerate}
\item Defina la ley de Avogadro, de formulas y constantes.
\item Defina la ley de Dalton, y Amagat, dar fórmulas.
\item De la fórmula de la cantidad de calor suministrada por un gas,
  indique los componentes de la misma. Explique calor específico a
  presión y volumen constante.
\item De un ejemplo con un problema de aplicación. Elijan un problema.
\end{enumerate}

\actividad{ Ejercicios.}
\begin{enumerate}
\item Determinar el \emph{volumen molar} que posee un gas de la
  atmósfera a $27^\circ C$ y $p=1Atm= 10330\frac{kg}{m^2}g$.
\item Calcular el \emph{número de moles} de una masa gaseosa que
  ocupa un volumen de $v=0,4m^3$ a la presión $p=2,5atm$ y a la
  temperatura $t=-3^\circ C$
\item Una mezcla de nitrógeno y oxigeno que a la temperatura
  $t=27^\circ C$ y a la presión de $p=1atm$, pesa $100Kg$, contiene
  un $40\%$ en peso de nitrógeno. Calcular las presiones parciales de cada gas.
\item Una mezcla de $48kg$ de oxígeno, 42kg de nitrógeno y 10kg de
  hidrógeno se encuentra a la temperatura de $47^\circ$ y a la
  presión $p=1,5atm$. Calcular el número de moles y el volumen de la mezcla.
\item ¿Cuál es el volumen molar de un gas de la atmósfera a
  27$^\circ$C y una presión de 1 atmósfera?
\item ¿Cuántos moles de gas hay en un volumen de 0.4 m$^3$ a una
  presión de 2.5 atmósferas y una temperatura de -3$^\circ$C?
\item Una mezcla de 48 kg de oxígeno, 52 kg de nitrógeno y 15 kg de
  hidrógeno se encuentra a una temperatura de 37$^\circ$C y una
  presión de 2 atmósferas. Calcular el número de moles y el volumen
  de la mezcla.
\item Si un gas ocupa un volumen de 2.5 L a una temperatura de
  25$^\circ$C y una presión de 1 atmósfera, ¿cuántos moles de gas hay?
\item ¿Cuál es la presión parcial del nitrógeno en una mezcla de
  gases que contiene 0.5 moles de nitrógeno y 1.5 moles de oxígeno si
  la presión total de la mezcla es de 2 atmósferas?
\item ¿Cuál es la presión de un gas que ocupa un volumen de 5 L a una
  temperatura de 27$^\circ$C si se sabe que contiene 0.1 moles de gas?
\item Si un gas tiene una presión de 3 atmósferas, un volumen de 2 L
  y una temperatura de 273 K, ¿cuál es el número de moles de gas?
\item ¿Cuál es la presión de un gas que ocupa un volumen de 10 L a
  una temperatura de 27°C si se sabe que contiene 0.1 moles de gas?
\item Si un gas tiene una presión de 1 atmósfera, un volumen de 10 L
  y una temperatura de 300 K, ¿cuál es la masa del gas si su masa
  molar es de 32 g/mol?
\item Un gas se encuentra en un recipiente con una presión de 2
  atmósferas y una temperatura de 25°C. Si se le añade más gas al
  recipiente hasta que la presión aumenta a 3 atmósferas y la
  temperatura se mantiene constante, ¿cuál es la razón entre el
  volumen final y el volumen inicial del gas?
\end{enumerate}
\begin{table}[h]
\centering
\begin{tabular}{|c|c|}
  \hline
  Fórmula & Descripción \\ \hline
  $PV = nRT$ & Ley de los gases ideales \\ \hline
  $N = \frac{P V}{R T}$ & Número de moles \\ \hline
  $V_m = \frac{V}{n}$ & Volumen molar \\ \hline
  $R = 0.0821 \frac{\textrm{L atm}}{\textrm{mol K}}$ & Constante
  universal de los gases \\ \hline
  $N_A = 6.0221 \times 10^{23}
  \frac{\textrm{moléculas}}{\textrm{mol}}$ & Número de Avogadro \\ \hline
  $M = \frac{m}{n}$ & Masa molar \\ \hline
\end{tabular}
\caption{Fórmulas y constantes para la Ley de Avogadro.}
\label{tabla:avogadro}
\end{table}

\end{document}

